% ==============================================================================
% sections/methodology.tex
% Created: 2026-05-08
% ==============================================================================

\subsection{Baseline model}\label{subsec:baseline}

\subsubsection*{Setup and notation}

Let $h_{it} \in \{0,1\}$ denote the true employment status of individual $i$ at
quarter $t$, where $h_{it} = 1$ indicates employment. We assume that the true
process follows a first-order, time-homogeneous Markov chain with transition
probabilities
\[
  \theta_0 = \Pr(h_{t}=1 \mid h_{t-1}=0), \qquad
  \theta_1 = \Pr(h_{t}=0 \mid h_{t-1}=1),
\]
where we suppress the individual subscript for clarity. The stationary
(long-run) employment rate is $\alpha = \theta_0 / (\theta_0 + \theta_1)$.
We consider both a stationarity-restricted specification in which the initial
probability $\Pr(h_1=1) = \alpha$ and an unrestricted specification in which
$\Pr(h_1=1)$ is estimated freely.

Survey respondents report $s_{it} \in \{0,1\}$, their observed employment
status, which may differ from $h_{it}$ due to measurement error. We assume
a symmetric, time-invariant misclassification structure:
\[
  \pi = \Pr(s_{t} = 1 \mid h_{t} = 0) = \Pr(s_{t} = 0 \mid h_{t} = 1),
\]
so that employment and non-employment are equally likely to be misclassified,
and errors are independent across waves conditional on the true state. The
full technical derivation of the likelihood and its properties, including
the identification analysis, is given in Appendix~\ref{app:baseline}.

\subsubsection*{Observed histories and marginal likelihood}

For each individual observed in three consecutive waves, the observed
employment history is the triple $(s_1, s_2, s_3) \in \{0,1\}^3$, yielding
eight possible observed histories $\mathcal{H} = \{000, 001, \ldots, 111\}$.
The marginal likelihood for an individual with observed history $\mathbf{s}$
is obtained by summing the joint likelihood of $(\mathbf{s}, \mathbf{h})$
over all eight latent histories $\mathbf{h} \in \mathcal{H}$:
\begin{equation}
  L_i(\theta_0, \theta_1, \pi) =
  \sum_{\mathbf{h} \in \mathcal{H}}
    \Pr(\mathbf{s}_i \mid \mathbf{h}; \pi)
    \Pr(\mathbf{h}; \theta_0, \theta_1),
  \label{eq:marginal_lik}
\end{equation}
where $\Pr(\mathbf{s} \mid \mathbf{h}; \pi) = \prod_{t=1}^{3} \Pr(s_t \mid h_t; \pi)$
exploits the conditional independence of errors across waves, and
$\Pr(\mathbf{h}; \theta_0, \theta_1)$ is the Markov chain probability of the
latent history. The sample log-likelihood is $\ell = \sum_i \log L_i$.

The marginal likelihood \eqref{eq:marginal_lik} is a finite mixture: at each
parameter draw, the data assign responsibility for each observation to the
eight latent histories in proportion to their joint likelihood. This mixture
structure is what makes the EM algorithm well-suited to this problem.

\subsubsection*{Identification}

The three-wave observed employment distribution contains eight cell probabilities
subject to one adding-up constraint, yielding seven free moments. The baseline
model has three free parameters: $\theta_0$, $\theta_1$, and $\pi$ (and,
in the free initial-probability specification, a fourth parameter for
$\Pr(h_1=1)$). Since seven moments exceed three (or four) free parameters,
the model is over-identified when the true process satisfies the first-order
Markov property.

The key identification insight is that the eight observed cell probabilities
encode the Markov property as a testable restriction. Specifically, the
conditional independence $s_1 \perp s_3 \mid s_2$ holds for any first-order
Markov process on \emph{observed} status---but it fails when the observed
process is a noisy version of a latent Markov chain. In that case, knowing
$s_1$ provides information about $h_2$ (because $h_2$ is correlated with $h_1$
through the latent dynamics), and hence information about $s_3$, even
conditional on $s_2$. The magnitude of this additional dependence identifies
$\pi$: if misclassification is zero, the observed process is itself Markov and
satisfies the restriction; as $\pi$ increases, the Markov violation grows.
With three waves, the restriction is testable and $\pi$ is non-parametrically
identified up to the auxiliary assumption of a symmetric error structure.

\subsubsection*{EM algorithm}

The EM algorithm maximises the complete-data log-likelihood---the
log-likelihood of the full $(s, h)$ data---by iterating between two steps.
The E-step computes, for each individual and each latent history
$\mathbf{h}$, the posterior responsibility:
\begin{equation}
  \gamma_{i\mathbf{h}}^{(k)} =
  \frac{\Pr(\mathbf{s}_i \mid \mathbf{h}; \pi^{(k)}) \Pr(\mathbf{h}; \theta_0^{(k)}, \theta_1^{(k)})}
       {L_i(\theta_0^{(k)}, \theta_1^{(k)}, \pi^{(k)})},
  \label{eq:estep}
\end{equation}
where $k$ indexes the iteration. The M-step updates the parameters by
maximising the expected complete-data log-likelihood:
\begin{align}
  \theta_0^{(k+1)} &= \frac{\sum_i \sum_{\mathbf{h}: h_1=0, h_2=1} \gamma_{i\mathbf{h}}^{(k)}
                            + \sum_i \sum_{\mathbf{h}: h_2=0, h_3=1} \gamma_{i\mathbf{h}}^{(k)}}
                           {\sum_i \sum_{\mathbf{h}: h_1=0} \gamma_{i\mathbf{h}}^{(k)}
                            + \sum_i \sum_{\mathbf{h}: h_2=0} \gamma_{i\mathbf{h}}^{(k)}}, \label{eq:mstep_theta0}\\
  \theta_1^{(k+1)} &= \frac{\sum_i \sum_{\mathbf{h}: h_1=1, h_2=0} \gamma_{i\mathbf{h}}^{(k)}
                            + \sum_i \sum_{\mathbf{h}: h_2=1, h_3=0} \gamma_{i\mathbf{h}}^{(k)}}
                           {\sum_i \sum_{\mathbf{h}: h_1=1} \gamma_{i\mathbf{h}}^{(k)}
                            + \sum_i \sum_{\mathbf{h}: h_2=1} \gamma_{i\mathbf{h}}^{(k)}}, \label{eq:mstep_theta1}\\
  \pi^{(k+1)}     &= \frac{\sum_i \sum_{\mathbf{h}} \gamma_{i\mathbf{h}}^{(k)}
                           \sum_{t=1}^{3} \mathbf{1}[s_{it} \ne h_t]}
                          {3\, \sum_i 1}. \label{eq:mstep_pi}
\end{align}
The M-step updates \eqref{eq:mstep_theta0}--\eqref{eq:mstep_pi} have a
transparent interpretation: $\theta_0$ is the responsibility-weighted empirical
fraction of periods in which a truly non-employed worker becomes employed,
$\theta_1$ is the analogous exit fraction, and $\pi$ is the
responsibility-weighted empirical misclassification rate. Each update is
in closed form, so the algorithm requires no numerical optimisation and
is guaranteed to converge to a local maximum of the observed-data
log-likelihood \citep{dempster1977}.

For the covariate specifications, the reported standard errors use an
individual-level survey-weighted sandwich covariance matrix. We obtain the
observed-information ``bread'' by differentiating the analytic observed-data
score and use the outer product of weighted individual scores as the ``meat.''
Standard errors for risk-set weighted rates, flows, and discrete effects are
then obtained by the delta method, including both parameter uncertainty and
direct sampling variation in the weighted empirical functionals. These are
linearisation approximations and do not incorporate survey strata or PSU
clustering, which are unavailable in the estimation extract. A nonparametric
individual-level bootstrap that resamples complete three-wave histories and
retains their survey weights and covariates is retained as a robustness option.

\subsection{Extensions}\label{subsec:extensions}

The extensions below each generalise one aspect of the baseline model. All
retain the core EM structure; the main difference is in the M-step update
for the extended parameter set.

\subsubsection{Observable worker heterogeneity}\label{subsubsec:covariates}

We allow transition probabilities to depend on individual and origin-wave characteristics
through probit link functions. Age, education, gender, and race enter both
transition equations. Current contract type is defined only for employed
respondents, so it enters the employment-persistence (exit-risk) equation only
and is excluded from the job-entry equation. An indicator for informal-sector
employment (including agriculture) is treated in the same way:
\[
  \theta_{0it} = \Phi(X_{0it} \beta_0), \qquad
  \theta_{1it} = \Phi(X_{1it} \beta_1),
\]
where $X_{1it}$ additionally contains contract type and the informal-sector
indicator measured at the origin wave (wave 1 for the first transition and
wave 2 for the second). These characteristics may therefore change between
transitions. We do not split employment into formal and informal latent states.
Because transition probabilities vary over time in Set 3, the main comparison
estimates the initial employment probability freely rather than imposing a
time-homogeneous stationary distribution. The function $\Phi(\cdot)$ is the
standard normal CDF. The misclassification parameter
$\pi$ is held common across individuals in the main specification; a variant
allows $\pi$ to vary with observable reliability indicators (see
Section~\ref{subsubsec:inconsistency}). Since the M-step for $\beta_0$ and
$\beta_1$ no longer has a closed form, we use the Generalised EM (GEM)
algorithm of \citet{dempster1977}. Under stationarity the initial-state term
couples the two coefficient vectors, so we update them jointly using limited
BFGS steps on the full expected complete-data log-likelihood. Backtracking is
applied until both that objective and the observed-data likelihood are
non-decreasing.

\subsubsection{Finite mixture model for unobserved heterogeneity}\label{subsubsec:fmm}

Unobserved worker heterogeneity---workers who persistently have very high or
very low transition probabilities---can mimic the Markov violations created by
misclassification: a persistently high-exit worker looks, in the observed data,
like someone who repeatedly enters and exits employment within the observation
window. We disentangle these mechanisms through a two-type finite mixture model
(FMM):
\[
  L_i(\phi, \theta^A, \theta^B, \pi) =
  \phi\, L_i(\theta^A, \pi) + (1-\phi)\, L_i(\theta^B, \pi),
\]
where type $A$ has parameters $(\theta_0^A, \theta_1^A)$, type $B$ has
parameters $(\theta_0^B, \theta_1^B)$, and $\phi$ is the mixing proportion.
The EM algorithm for the FMM introduces a second layer of latent variables
(the type indicator) and iterates the E-step responsibilities jointly over
latent states and types. The misclassification parameter $\pi$ is identified
separately from the type heterogeneity because Markov violations attributable
to types can be attributed to heterogeneity in $(\theta^A, \theta^B, \phi)$, while
additional violations residual to the type structure are attributed to $\pi$.

\subsubsection{Inconsistency-augmented misclassification}\label{subsubsec:inconsistency}

If misclassification is driven by survey response error, respondents who make
other types of survey errors should also have higher misclassification rates.
We test this prediction by allowing $\pi$ to depend on an indicator of survey
unreliability. Specifically, we define:
\begin{itemize}
  \item $A_i = 1$ if individual $i$ has an age inconsistency in any wave pair;
  \item $E_i = 1$ if individual $i$ has an education inconsistency in any wave pair.
\end{itemize}
We then estimate the generalised model:
\[
  \pi_i = \mathrm{logit}^{-1}(\beta_0 + \beta_1 A_i + \beta_2 E_i),
\]
where $\mathrm{logit}^{-1}(x) = e^x/(1+e^x)$ ensures that $\pi_i \in (0,1)$.
The identifying variation comes from whether the inconsistency is attributed
to the first, second, or third wave of the triad. We use a wave-attribution
rule based on the sign of the change in the inconsistent variable to assign
each inconsistency to the wave most likely responsible, and we then condition
$\pi_i$ on the presence of a second-wave attribution (since a first-wave error
in the first-wave status and a third-wave error in the third-wave status
would be absorbed into the initial and terminal probabilities). The details of
the attribution rule are described in Appendix~\ref{app:baseline}.

\subsubsection{AR(2) extension}\label{subsubsec:ar2}

The baseline model assumes that the latent true employment process is
first-order Markov: conditional on last quarter's true status, the probability
of this quarter's status is independent of all earlier history. This rules out
duration dependence---the possibility that the probability of finding a job
increases or decreases with the length of the current non-employment spell.
We relax this assumption by extending the model to a second-order Markov
specification using four consecutive waves.

In the AR(2) model, the latent transition probability depends on the current
state and the state one period earlier:
\begin{align*}
  p_{00} &= \Pr(h_t=1 \mid h_{t-1}=0, h_{t-2}=0), &
  p_{01} &= \Pr(h_t=1 \mid h_{t-1}=0, h_{t-2}=1), \\
  p_{10} &= \Pr(h_t=0 \mid h_{t-1}=1, h_{t-2}=0), &
  p_{11} &= \Pr(h_t=0 \mid h_{t-1}=1, h_{t-2}=1),
\end{align*}
where the first subscript denotes the state at $t-1$ and the second the state
at $t-2$. This specification accommodates negative duration dependence in
non-employment (if $p_{00} > p_{01}$, a longer non-employment spell reduces
the probability of entry) and positive duration dependence in employment. The
four-wave observed histories, of which there are 16, support identification of
these four additional transition probabilities along with the misclassification
parameter. The full derivation is in Appendix~\ref{app:ar2}.

\subsubsection{Tenure contamination model}\label{subsec:tenure}

The baseline model and its extensions above identify misclassification from
Markov violations in the employment status sequence. The tenure contamination
model provides an independent identification channel using job tenure and
non-employment duration data.

The key observation is that if an employment observation is misclassified---say,
a truly non-employed worker is recorded as employed at wave $t$---the reported
job tenure at wave $t+1$ should be inconsistent with normal tenure accumulation.
Specifically, if the worker was truly non-employed at $t$ and is truly employed
at $t+1$, their tenure at $t+1$ should be short (consistent with a new job).
But if the worker is also truly employed at $t-1$, their tenure at $t-1$ should
be long (consistent with a continuing job), and the jump in tenure between
$t-1$ and $t+1$ (or the failure of tenure to accumulate between $t$ and $t+1$)
signals a ``contaminated'' intermediate observation.

We formalise this intuition as follows. Let $d_{it}$ denote the duration
variable (tenure if employed, time since last job if non-employed) at wave $t$.
We model the joint distribution of $(s_{i,t-1}, d_{i,t-1}, s_{it}, d_{it},
s_{i,t+1}, d_{i,t+1})$ as a function of the underlying latent employment
sequence $\mathbf{h}$ and an additional contamination parameter $\varepsilon$:
the probability that a duration observation is ``contaminated,'' meaning it is
drawn from the distribution of the wrong employment state. The contamination
probability is interpreted as a form of tenure recall error that is correlated
with employment state misclassification. Jointly identifying $\pi$ and
$\varepsilon$ from the observed employment--tenure sequence requires an
assumption about the marginal distribution of tenure; we model tenure as
exponential with rate parameters $(\lambda_g, \lambda_d)$ for employed and
non-employed workers respectively, and fit the model by EM. The full derivation
is in Appendix~\ref{app:tenure}.

\subsection{Matching and sample construction}\label{subsec:matching}

Each three-wave triad is constructed by linking panel observations using
household and individual identifiers. We consider a hierarchy of matching
strictness, from the baseline matching---which requires agreement on household
identifiers only---to stricter matching algorithms that additionally require
agreement on gender, race, and certain time-invariant characteristics. Stricter
matching reduces the sample size (by discarding potential false matches) but
should also reduce the incidence of spurious transitions arising from
mismatched observations rather than genuine employment changes. As discussed
in Section~\ref{sec:alternatives}, we use this variation as a design-based
robustness check: if our estimated misclassification $\pi$ reflects genuine
survey error rather than modelling artefact, stricter matching should reduce
but not eliminate $\pi$.
